\documentclass[10pt, oneside]{article}   	% use "amsart" instead of "article" for AMSLaTeX format
\usepackage{geometry}                		% See geometry.pdf to learn the layout options. There are lots.
\geometry{letterpaper}                   		% ... or a4paper or a5paper or ... 
%\geometry{landscape}                		% Activate for rotated page geometry
%\usepackage[parfill]{parskip}    		% Activate to begin paragraphs with an empty line rather than an indent
\usepackage{graphicx}				% Use pdf, png, jpg, or eps§ with pdflatex; use eps in DVI mode
								% TeX will automatically convert eps --> pdf in pdflatex		
\usepackage{amssymb}

\usepackage[T2A]{fontenc}
\usepackage[utf8]{inputenc}
\usepackage[russian]{babel}

\usepackage{tipa}


\usepackage{caption}


\usepackage{tikz}

\newcommand{\liaison}{%
  \tikz[baseline=-8pt] \draw (-0.1,0) to[bend left=30] (0.1,0);%
}

%SetFonts

%SetFonts

\title{Le française A1}
\author{Anton Yesenin}
\date{}							% Activate to display a given date or no date

\begin{document}
\maketitle
%\section{}
%\subsection{}
\pagebreak

\section{Урок 1}

%ʃ \textesh
%ʒ \textyogh
%ɛ \textepsilon

\subsection{Фонетическое упражнение 1}
[pa~--- ba | fa~--- va | ta~--- da | sa~--- za | {\textesh}a~--- {\textyogh}a | ka~--- ga | ma~--- na]
\subsection{Фонетическое упражнение 2}
[p{\textepsilon}~--- b{\textepsilon} | f{\textepsilon}~--- v{\textepsilon} | t{\textepsilon}~--- d{\textepsilon} | s{\textepsilon}~--- z{\textepsilon} | {\textesh}{\textepsilon}~--- {\textyogh}{\textepsilon} | k{\textepsilon}~--- g{\textepsilon} | m{\textepsilon}~--- n{\textepsilon}]
\subsection{Фонетическое упражнение 3}
[pi~--- bi | fi~--- vi | ti~--- di | si~--- zi | {\textesh}i~--- {\textyogh}i | ki~--- gi | mi~--- ni]
\subsection{Фонетическое упражнение 4}
[pal~--- bal | f{\textepsilon}l~--- v{\textepsilon}l | til~--- dil | sal~--- zal | {\textesh}{\textepsilon}l~--- {\textyogh}{\textepsilon}l | kil~--- gil | mal~--- nal]
\subsection{Фонетическое упражнение 5}
[pl{\textepsilon}~--- bl{\textepsilon} | fla~--- vla | tl{\textepsilon}~--- dl{\textepsilon} | sla~--- zla | {\textesh}l{\textepsilon}~--- {\textyogh}l{\textepsilon} | kla~--- gla | ml{\textepsilon}~--- nl{\textepsilon}]
\subsection{Фонетическое упражнение 6}
[pa:r~--- ba:r | fa:r~--- va:r | {\textesh}a:r~--- {\textyogh}a:r | ka:r~--- ga:r | ma:r~--- na:r~--- ra:r]
\subsection{Фонетическое упражнение 7}
[pra~--- bra | fra~--- vra | tra~--- dra | sra~--- zra | {\textesh}ra~--- {\textyogh}ra | kra~--- gra | mra~--- nra | lra~--- ra]
\subsection{Фонетическое упражнение 8}
\subsection{Фонетическое упражнение 9}
\subsection{Фонетическое упражнение 10}

\subsection{Упражнение в чтении A}
\subsubsection{1}
$\square$...$\square$e $\to$ [$\square$...$\square$]\\
\slshape salad{\bfseries e}, parad{\bfseries e}, artist{\bfseries e}, madam{\bfseries e}, mani{\bfseries e}, pirat{\bfseries e}, Mari{\bfseries e}, Itali{\bfseries e}, Albani{\bfseries e}, vals{\bfseries e}, litr{\bfseries e}, ministr{\bfseries e}, titan{\bfseries e}, dam{\bfseries e}, dram{\bfseries e}
\upshape

\subsubsection{2}
a = à = â = [a]\\
\slshape p{\bfseries a}n{\bfseries a}m {\bfseries a}, {\bfseries à}, l{\bfseries à}, {\bfseries â}me, p{\bfseries â}te, r{\bfseries â}le, {\bfseries â}ne, B{\bfseries a}rb{\bfseries a}r{\bfseries a}
\upshape

\subsubsection{3}
$\square$...$\square$s $\to$ [$\square$...$\square$]\\
\slshape Pari{\bfseries s}, pa{\bfseries s}, paradi{\bfseries s}, canari{\bfseries s}, pi{\bfseries s}, panama{\bfseries s}, pli{\bfseries s}, di{\bfseries s}, li{\bfseries s},  ri{\bfseries s}
\upshape

\subsubsection{4}
bananes, articles, fatales, palmes, pirates, salades, dates, banales, parades, artistes
\subsubsection{5}
salle, balle, Vannes, Anne, panne, dalle, patte
\subsubsection{6}
belle, reste, taverne, serbe, adresse, Berne, Rennes, Ardennes, Bernardette, Annette
\subsubsection{7}
race, cidre, pacifiste, citadelle, milice, carpe, acte, crabes, cabine, classe, Cannes, Canada, câble, carte, barricade, place, mascarade, crabe, caravane, pacte, caprice, caste, spectacle
\subsubsection{8}
certificate, apparat, primat, plat, climat

\subsection{Упражнение в чтении B}
ma, datte, ta, sa, ci, la, palme, primitive, pis, part, fatale, victime, dette, canne, banale, titane, table, âpre, admire, article, pâte, trace, cercle, mât, terre, rat, telle, citerne, ami, carte, traverse, tapis, fictive, certes, ici, parties, capable, visites, apparat, bras, âme
\pagebreak

\section{Урок 2}

\section{Урок 3}
Elle{\liaison}habite{\liaison}à{\liaison}Paris.\\
Il va{\liaison}à{\liaison}Arles.
\end{document}  